\section*{Introduction}
\addcontentsline{toc}{section}{\protect\numberline{}Introduction}
\label{sec:introduction}

Le développement de fusées sondes et de lanceurs expérimentaux constitue aujourd’hui un terrain privilégié d’innovation
technologique, notamment dans le cadre de projets étudiants et associatifs. Ces systèmes, soumis à de fortes contraintes de
masse, de fiabilité et de sécurité, nécessitent des solutions de conception à la fois simples, robustes et économes en
ressources. Parmi les problématiques clés rencontrées lors de la conception de lanceurs multi-étages figure la séparation
d’étages, étape critique du vol dont la fiabilité conditionne directement le succès de la mission.\\

Les systèmes de séparation conventionnels reposent majoritairement sur des dispositifs pyrotechniques ou des mécanismes à forte
complexité mécanique. Bien que largement éprouvées dans le domaine spatial industriel, ces solutions présentent plusieurs
inconvénients dans le cadre de projets expérimentaux de petite échelle : complexité de mise en œuvre, contraintes
réglementaires, coûts élevés et risques accrus lors des phases de test. Dans ce contexte, l’étude de solutions alternatives de
séparation apparaît comme une piste pertinente, en particulier pour des lanceurs suborbitaux de type fusée sonde.\\

L’aérofreinage, habituellement utilisé pour la stabilisation ou la décélération d’un véhicule en phase atmosphérique, peut
également être envisagé comme un moyen de générer un effort différentiel suffisant pour provoquer la séparation contrôlée de deux
étages. Toutefois, l’intégration d’un tel principe soulève plusieurs questions techniques majeures : capacité du dispositif à
générer les efforts requis, interaction avec l’écoulement aérodynamique global, tenue structurelle des éléments déployables et
compatibilité avec l’architecture générale du lanceur.\\

C’est dans ce cadre que s’inscrit le présent travail, réalisé au titre du \acrfull{PMI} de l’\acrfull{IPSA} et intitulé :

\begin{quote}
    \centering
    \textit{\og{}Perspective d’expérience d’un système de séparation d’étages de fusée par aérofreinage.\fg{}}
\end{quote}

\noindent Cette étude prend appui sur le projet de l'association étudiante AéroIPSA\cite{AeroIPSA} : SP-02, qui constitue le
support expérimental et le cadre d’intégration du système étudié. Il convient toutefois de souligner que le projet SP-02 n’est pas
l’objet de ce rapport : il fournit un environnement géométrique, structurel et aérodynamique réaliste, nécessaire à la conduite
des analyses menées dans le cadre du \acrshort{PMI}.\\

L’objectif de ce travail est d’évaluer la faisabilité technique d’un système de séparation d’étages reposant sur l’aérofreinage,
à travers une démarche combinant étude théorique, modélisation numérique et conception préliminaire. L’analyse porte en
particulier sur la capacité du dispositif à générer un effort aérodynamique exploitable, sur son intégration mécanique au sein
du lanceur et sur la cohérence physique des hypothèses retenues. Les systèmes d’aérofreinage et de séparation sont étudiés
indépendamment de toute solution pyrotechnique.\\

Le périmètre du \acrshort{PMI} est volontairement limité à l’étude, la conception et l’analyse des sous-systèmes directement liés
à la séparation par aérofreinage. Les autres éléments du lanceur (propulsion, récupération, avionique globale) ne sont abordés que
dans la mesure où ils influencent les conditions aux limites des simulations aérodynamiques et structurelles. Une présentation
synthétique de l’architecture générale du lanceur est ainsi proposée afin de contextualiser les choix de modélisation retenus,
sans constituer une description exhaustive du projet SP-02.\\

Le rapport est organisé de la manière suivante. La première partie présente le contexte technique et le cadre de l’étude,
incluant une description synthétique du lanceur servant de système hôte. La seconde partie détaille la méthodologie adoptée, les
hypothèses de modélisation et les outils de simulation utilisés. La troisième partie est consacrée à la conception du système de
séparation et des aérofreins. Les résultats issus des simulations aérodynamiques et structurelles sont ensuite présentés et
analysés de manière critique. Enfin, une dernière partie aborde les perspectives de validation expérimentale et les limites de
l’étude, avant de conclure sur les apports et les prolongements possibles de ce travail.

\newpage
